\documentclass[12pt]{article}
\usepackage[utf8]{inputenc}
\usepackage[T1]{fontenc}
\usepackage{hyperref}
\title{Zayse Language Resources for K-12 Curriculum Development}
\author{Compiled by Claude Sonnet 5.5 (claude-sonnet-5-5) for Andr\'as Kornai}
\date{}
\begin{document}
\maketitle

\section{Introduction}
Zayse (Zaysete; ISO 639-3 zay; Glottocode zays1235) is an Omotic language of southern Ethiopia. It is used as a first language in its community and is taught as a subject in some schools. Speaker counts in the sources vary widely, from about 10,000 to 34,000, probably because of differing ethnic and dialect boundaries. Little pedagogical material turned up in the searches. What exists is mostly linguistic documentation and religious audio.

\section{Primary Educational Resources}
No primers, textbooks or school curricula for Zayse turned up in the searches. A curriculum would have to be built from the documentation listed below, ideally with community and regional education bureau involvement.

\section{Linguistic and Reference Materials}
\begin{itemize}
\item Documentation and Grammatical Description of Zaysete, a PhD dissertation at Addis Ababa University. It covers phonology, morphology and syntax, and is reported to include annotated audio and video recordings. \url{https://etd.aau.edu.et/items/b757684f-86c0-4142-b6d3-ce2700435273}
\item Full record of the same dissertation. \url{https://etd.aau.edu.et/items/b757684f-86c0-4142-b6d3-ce2700435273/full}
\item LINGUIST List announcement (issue 30.4667) describing the dissertation. \url{https://linguistlist.org/issues/30/4667/}
\item Wikipedia article on Zayse-Zergulla, which lists further scholarly references. \url{https://wikipedia.com/wiki/Zayse}
\item Endangered Languages Project entry. \url{https://endangeredlanguages.com/node/91890}
\end{itemize}

\section{Religious and Community Resources}
\begin{itemize}
\item Global Recordings Network, Zayse recordings. \url{https://globalrecordings.net/en/language/zay}
\item Joshua Project, Zayse people group profile. \url{https://www.joshuaproject.net/people_groups/16096}
\item Joshua Project, Zayse language page. \url{https://www.joshuaproject.net/index.php/languages/zay}
\end{itemize}

\section{Digital and Audio Resources}
\begin{itemize}
\item Wikitongues language profile. \url{https://wikitongues.org/languages/zay/}
\item The audio and video recordings attached to the dissertation (see above) are the main digital corpus found.
\end{itemize}

\section{Curriculum Development Context}
Zayse is a stable language that children still acquire, so mother-tongue instruction is plausible. The searches found no standardized teaching materials. Priorities would be an agreed orthography, primers drawn from the documented texts, and teacher training. Local authorities and speakers should be consulted. Anything not listed here has not been verified.
\end{document}